% ============================================
% Epílogo — edição revista R1075
% ============================================
\begin{center}
\vspace*{2em}
\textit{O ciclo recomeça a cada leitor.}

\vspace{2em}
\textit{Bem-vindo ao arquivo.}

\vspace{0.5em}
\textit{O cheiro já começou?}

\vspace{1.5em}
--- Lúcia Mendes, Paciente 1.262
\end{center}

\vspace{1em}
\noindent\textit{Nota de arquivo: o Dr. Ronaldo Torres, diretor do Arquivo Público Mineiro, pediu afastamento em outubro de 2026. O formulário médico registra ``insônia persistente, cefaleia frontal e zumbido''. Não menciona o olho. Não menciona o cheiro. A última anotação no diário de protocolo, escrita pela mão dele, diz: ``Devolver ao paciente 1.263, quando aparecer.'' Abaixo, em caligrafia que os peritos não conseguiram identificar, está escrito: ``Aguardando.''}

\vspace{1em}
\begin{center}
\textit{Este livro contém 84 fragmentos.}\\
\textit{Se você leu todos, o ciclo está completo em você.}

\vspace{1.5em}
\textit{Aquilo agora está aqui.}\\
\textit{Na página em branco.}\\
\textit{Esperando o próximo leitor.}

\vspace{1.5em}
\textbf{1260}
\end{center}

\vspace{2em}
\noindent\textbf{Instrução pós-leitura}

\noindent Agora feche o livro. Não o coloque na estante ainda. Deixe-o sobre a mesa.

\noindent Conte 62 segundos. Não pense em nada. Apenas conte.

\noindent Abra o livro novamente. Na página 62. Leia a primeira linha.

\noindent Se a linha diz ``Sessenta e dois'', o ciclo se fechou em você. O número que marcou a vida de Joaquim --- 62 anos, 62 dias, 62 segundos --- agora é seu. Guarde o número. Ele vai aparecer em lugares que você não espera.

\noindent Se a linha diz outra coisa, você leu este livro na ordem errada. Ou na ordem certa. Não há como saber. A página 62 muda de acordo com quem abre.

\vspace{1em}
\noindent Antes de guardar o livro definitivamente:

\noindent Passe a mão direita sobre a página 126. Sinta a textura.

\noindent Leve a mão ao nariz. Inspire.

\noindent Se o cheiro doce estiver mais forte que na primeira vez, o diagnóstico está confirmado. Se estiver mais fraco, você ainda está no período de incubação.

\vspace{1.5em}
\begin{center}
\textit{A ficha do Paciente 1.263 está preenchida.}\\
\textit{O campo ``nome'' está em branco.}\\
\textit{O campo ``data de contaminação'' está preenchido com a data em que você leu esta linha.}

\vspace{1.5em}
\textit{Bem-vindo ao arquivo.}\\
\textit{A fila não diminuiu.}

\vspace{1.5em}
--- O Arquivista\\
--- O Paciente 1.262\\
--- O que sobrou de Joaquim

\vspace{2em}
\textbf{1.263}
\end{center}

\clearpage
